\chapter{Computer Languages}
\label{chap:computer_languages}

\begin{tcolorbox}[enhanced, colback=boxnavybg, colframe=brandnavy, boxrule=1.2pt, arc=4pt, title={\Large\bfseries\sffamily Unit 1 Overview \& Learning Objectives}]
\sffamily\normalsize
This chapter introduces the formal systems used to communicate instructions to digital computers. By the conclusion of this unit, students will be able to:
\begin{itemize}[leftmargin=1.2em, itemsep=2pt, topsep=3pt]
    \item Understand the structural hierarchy of computer languages from machine code to high-level abstractions.
    \item Differentiate between compilers, interpreters, and assemblers in terms of execution mechanics and error diagnostics.
    \item Trace the complete Program Development Cycle from problem specification to long-term software maintenance.
    \item Follow the step-by-step compilation and execution pipeline through preprocessing, compilation, assembly, linking, and loading.
    \item Analyze the internal CPU instruction cycle (Fetch, Decode, Execute, Store).
\end{itemize}
\end{tcolorbox}

\section{Orientation: Communicating with Computers}

Computer languages are structured, formal systems used to define unambiguous instructions that a digital computer can store, translate, and execute. Computers cannot understand conversational human speech directly; at the deepest hardware level, digital processors operate strictly through binary electrical logic states ($0$ and $1$). 

Over the history of computer science, programming representations evolved from tedious, machine-dependent binary sequences to human-readable high-level languages like C, Java, and Python. However, regardless of how advanced or expressive a programming language is, \textbf{the CPU ultimately executes only machine-language binary instructions}.

\begin{definition}[Core Principles of Computer Languages]
\begin{itemize}[leftmargin=1.2em, itemsep=3pt]
    \item \textbf{Binary Representation:} All computer instructions, memory addresses, and data values are internally represented as binary digits (\textbf{bits}), corresponding to physical electrical states (voltage thresholds).
    \item \textbf{Translation Requirement:} Because humans write in symbolic assembly or high-level code, translators (assemblers, compilers, or interpreters) are fundamentally required to produce executable machine code.
    \item \textbf{Machine Dependence:} Machine and assembly languages are tightly coupled to a specific processor family's Instruction Set Architecture (ISA), whereas high-level languages provide portability across diverse hardware architectures.
    \item \textbf{Syntax vs. Semantics:} A valid program must satisfy both \textbf{syntax} (grammatical character structure and statement format) and \textbf{semantics} (logical meaning and expected operational behavior).
\end{itemize}
\end{definition}

\section{Machine Language (Lowest Level)}

Machine language is the fundamental native programming language of any computer system. It consists entirely of strings of binary digits ($0$s and $1$s) that are directly decoded and executed by the processor's electronic logic gates without requiring any software translator.

\subsection{Instruction Format \& Structure}
A standard machine-language instruction is typically partitioned into two primary functional components:
\begin{equation}
\text{Machine Instruction} = [\;\text{\textbf{Opcode}}\;] \quad+\quad [\;\text{\textbf{Operand(s)}}\;]
\end{equation}
\begin{itemize}[leftmargin=1.2em]
    \item \textbf{Opcode (Operation Code):} A unique bit pattern that instructs the CPU's Control Unit which operation to perform (e.g., \texttt{ADD}, \texttt{SUB}, \texttt{LOAD}, \texttt{STORE}).
    \item \textbf{Operand:} Specifies the data or the physical memory addresses/registers on which the opcode must operate.
\end{itemize}

\begin{example}[Binary Machine Code Representation]
A typical 16-bit machine instruction might appear as:
\begin{center}
\texttt{10110000 01100001}
\end{center}
Here, the high byte (\texttt{10110000}) could represent the opcode command ``Move immediate byte into register AL'', while the low byte (\texttt{01100001}) represents the immediate hexadecimal data value \texttt{0x61} (ASCII character 'a'). The exact interpretation depends entirely on the processor's Instruction Set Architecture (e.g., x86 vs. ARM).
\end{example}

\subsection{Advantages \& Limitations of Machine Code}
\begin{itemize}[leftmargin=1.2em]
    \item \textbf{Advantages:} Provides direct execution without runtime translation overhead; delivers maximum hardware efficiency and granular control over processor resources.
    \item \textbf{Limitations:} Extremely difficult for humans to write, memorize, and debug; highly prone to catastrophic errors (altering a single bit changes the entire instruction); completely non-portable between different CPU architectures.
\end{itemize}

\section{Assembly Language (Symbolic Low Level)}

To overcome the immense difficulty of writing raw binary strings, computer scientists created \textbf{Assembly Language}. Assembly is a symbolic, low-level language that replaces binary opcodes with intuitive, short English-like abbreviations called \textbf{mnemonics}.

\subsection{Symbolic Instructions and Syntax}
An assembly language statement directly models the underlying processor instruction format:
\begin{lstlisting}[language={[x86masm]Assembler}]
START:  MOV R1, #5          ; Load immediate constant 5 into register R1
        ADD R2, R1, R1      ; Add R1 to R1 and place sum into register R2
        B   START           ; Unconditional branch back to label START
\end{lstlisting}
In this instruction sequence:
\begin{itemize}[leftmargin=1.2em]
    \item \texttt{MOV}, \texttt{ADD}, and \texttt{B} are \textbf{mnemonics} representing opcodes.
    \item \texttt{R1} and \texttt{R2} are internal \textbf{CPU registers}.
    \item \texttt{\#5} represents an \textbf{immediate operand}.
    \item \texttt{START:} is a \textbf{symbolic label} representing a memory address.
\end{itemize}

\subsection{One-to-One Relationship \& The Assembler}
Unlike high-level languages where a single statement can expand into dozens of machine operations, assembly language maintains an almost strict \textbf{one-to-one relationship}: one assembly statement translates into exactly one machine-code instruction.

Translating assembly language into machine-readable object code requires a dedicated translator known as an \textbf{Assembler}.

\section{High-Level Languages}

High-level programming languages (such as C, C++, Java, and Python) utilize natural English words, standard mathematical notation, and structured programming paradigms to hide processor-level hardware details from the software developer.

\subsection{Key Characteristics of High-Level Languages}
\begin{itemize}[leftmargin=1.2em]
    \item \textbf{Abstraction:} Calculations are expressed naturally (e.g., \texttt{grossPay = hours * rate}) without requiring the programmer to manually allocate registers, manage memory addresses, or issue bus read commands.
    \item \textbf{Portability:} A high-level program written on an Intel x86 machine can be recompiled or run on an ARM-based computer without major code modifications.
    \item \textbf{Structured Constructs:} Features such as loops (\texttt{for}, \texttt{while}), conditional branching (\texttt{if-else}), modular functions, and object-oriented classes allow developers to organize complex software systems into manageable units.
\end{itemize}

\section{Language Translators: Compiler, Interpreter, and Assembler}

Computers cannot execute high-level statements directly. Therefore, specialized system programs called \textbf{language translators} convert human-authored source text into executable machine instructions.

\subsection{The Compiler}
A \textbf{compiler} is a translator that converts an entire high-level source program into machine code or intermediate object code \textbf{all at once prior to program execution}.

\begin{definition}[Core Phases of a Modern Compiler]
\begin{enumerate}[leftmargin=1.2em]
    \item \textbf{Lexical Analysis (Scanner):} Reads source code characters and groups them into atomic units called \textbf{tokens} (identifiers, literals, operators, keywords).
    \item \textbf{Syntax Analysis (Parser):} Checks tokens against the formal grammatical rules of the language, constructing an Abstract Syntax Tree (AST). Rejects invalid grammar like \texttt{x = + 4;}.
    \item \textbf{Semantic Analysis:} Validates logical meaning, scope, and type compatibility (e.g., ensuring a string is not assigned to an integer).
    \item \textbf{Optimization:} Restructures intermediate code to run faster or consume less memory without changing program semantics.
    \item \textbf{Code Generation:} Emits final target machine code or assembly instructions for the host architecture.
\end{enumerate}
\end{definition}

\subsection{The Interpreter}
An \textbf{interpreter} translates and executes a high-level program \textbf{line-by-line during runtime}. Instead of generating a standalone executable binary beforehand, the interpreter reads a statement, parses it, executes the requested operation immediately, and then proceeds to the next statement.

\begin{center}
\small
\renewcommand{\arraystretch}{1.3}
\begin{tabularx}{\textwidth}{l X X}
\toprule
\textbf{Comparison Criterion} & \textbf{Compiler (e.g., C, C++)} & \textbf{Interpreter (e.g., Python, JS)} \\
\midrule
\textbf{Translation Timing} & Entire program translated before execution & Translated statement-by-statement during runtime \\
\textbf{Execution Speed} & Fast native execution; no translation overhead & Slower execution due to runtime translation \\
\textbf{Error Detection} & Reports all compile-time syntax errors at once & Halts immediately when the first error is reached \\
\textbf{Object Code} & Generates permanent object/executable files & Does not generate standalone machine code files \\
\textbf{Development Cycle} & Requires edit-compile-link-run cycles & Interactive testing and rapid prototyping \\
\bottomrule
\end{tabularx}
\end{center}

\subsection{The Assembler}
An \textbf{assembler} translates symbolic assembly code into binary machine code. Most modern assemblers use a \textbf{Two-Pass Architecture}:
\begin{itemize}[leftmargin=1.2em]
    \item \textbf{Pass 1 (Symbol Resolution):} Scans the entire file to find all label definitions and symbolic names, building a \textbf{Symbol Table} that maps each symbol to its assigned memory offset.
    \item \textbf{Pass 2 (Code Generation):} Scans the code a second time, replaces symbolic references with the actual memory addresses from the symbol table, and emits the binary object file.
\end{itemize}

\section{The Program Development Cycle}

Building reliable software requires following an organized engineering methodology known as the \textbf{Program Development Cycle}:

\begin{enumerate}[leftmargin=1.5em]
    \item \textbf{Problem Definition:} Clearly identifying the problem, defining expected outputs, inputs, operational constraints, and target users.
    \item \textbf{Problem Analysis:} Breaking down data requirements, validation checks (e.g., preventing negative pay hours), and boundary conditions.
    \item \textbf{Algorithm Design:} Developing a step-by-step, finite logical solution represented via pseudocode or standard flowcharts.
    \item \textbf{Coding:} Translating the algorithm into a concrete high-level programming language using proper variables, data types, and control structures.
    \item \textbf{Translation:} Using a compiler or interpreter to identify syntax errors and produce runnable machine code.
    \item \textbf{Testing \& Debugging:} Running the software with diverse test data (normal, boundary, and erroneous data) to detect and correct syntax, runtime, and logic errors.
    \item \textbf{Documentation:} Preparing user manuals and internal code comments explaining system architecture and maintenance guidelines.
    \item \textbf{Maintenance:} Ongoing updates to patch bugs, adapt to new operating systems, and add new capabilities over time.
\end{enumerate}

\section{The Compilation and Execution Process}

\subsection{The Complete Translation Pipeline}
When a programmer writes a program in a compiled language like C or C++, the source code undergoes several coordinated steps before executing on the processor:

\begin{center}
\fbox{\parbox{2.6cm}{\centering\textbf{Source Code}\\\texttt{main.c}}} $\longrightarrow$
\fbox{\parbox{2.6cm}{\centering\textbf{Preprocessor}\\\texttt{\#include} expansion}} $\longrightarrow$
\fbox{\parbox{2.6cm}{\centering\textbf{Compiler}\\Syntax/Semantic AST}} $\longrightarrow$
\fbox{\parbox{2.6cm}{\centering\textbf{Assembler}\\Object code \texttt{.obj}}} $\longrightarrow$
\fbox{\parbox{2.6cm}{\centering\textbf{Linker}\\Resolves \texttt{printf}}} $\longrightarrow$
\fbox{\parbox{2.6cm}{\centering\textbf{Loader}\\Loads to RAM}}
\end{center}

\begin{itemize}[leftmargin=1.2em]
    \item \textbf{Preprocessing:} Directives beginning with \texttt{\#} are handled (macro substitution and inclusion of header files).
    \item \textbf{Compilation:} Converts preprocessed code into assembly code.
    \item \textbf{Assembly:} Converts assembly code into relocatable object code files containing machine instructions and symbol tables.
    \item \textbf{Linking:} Merges independent object files with pre-compiled standard system libraries (e.g., standard input/output routines) to produce a standalone executable binary.
    \item \textbf{Loading:} The operating system's loader allocates address space in RAM, copies the executable segments into memory, and passes control to the CPU.
\end{itemize}

\subsection{The CPU Instruction Cycle}
Once instructions reside in primary memory, the central processor repeatedly performs the foundational \textbf{Instruction Cycle}:
\begin{enumerate}[leftmargin=1.5em]
    \item \textbf{Fetch:} The Control Unit retrieves the instruction from the memory address specified by the \textbf{Program Counter (PC)} register into the \textbf{Instruction Register (IR)}.
    \item \textbf{Decode:} The decoder circuit interprets the opcode bits to determine the required operation and operand locations.
    \item \textbf{Execute:} The Arithmetic Logic Unit (ALU) or control circuits execute the requested operation (such as adding two register values).
    \item \textbf{Store (Write-back):} The result of the computation is written back to a destination CPU register or memory location.
\end{enumerate}

\section{Chapter Review \& Key Takeaways}
\begin{tcolorbox}[enhanced, colback=boxgreenbg, colframe=boxgreenborder, boxrule=1pt, arc=4pt, title={\bfseries\sffamily Summary Checklist}]
\sffamily\small
\begin{itemize}[leftmargin=1.2em, itemsep=2pt]
    \item Machine language consists exclusively of binary digits ($0$ and $1$) and runs directly on the CPU without translation.
    \item Assembly language replaces binary opcodes with mnemonics (\texttt{MOV}, \texttt{ADD}) and uses an assembler to generate object code.
    \item Compilers translate the whole program at once before execution; interpreters evaluate code statement-by-statement at runtime.
    \item The compilation pipeline proceeds from Preprocessor $\rightarrow$ Compiler $\rightarrow$ Assembler $\rightarrow$ Linker $\rightarrow$ Loader.
    \item The CPU instruction execution cycle consists of four continuous phases: Fetch $\rightarrow$ Decode $\rightarrow$ Execute $\rightarrow$ Store.
\end{itemize}
\end{tcolorbox}
